% ------------------------------------------------------------------------
% ------------------------------------------------------------------------
% Modelo UNIAVAN para Trabalhos Acadêmicos baseado em Modelo UFSC de 
% Alisson Lopes Furlani utilizando a classe abntex2
%
% Autor: Luiz Fernando Marquez Arruda
%
% Olá caro autor, segue aqui algumas orientações a respeito de como utilizar este 
% template.
%
% Todas as configurações a respeito de nome, curso, orientador e afins devem ser
% feitas no arquivo information.tex que se encontra na pasta setup. Assim, de forma
% automática, os nomes corretos irão aparecer em sua capa, contra-capa e afins.
%
% Todas as imagens devem submetidas para a pasta images.
%
% A lista de símbolos, siglas e acrônimos também está localizada na pasta setup e
% consta com alguns exemplos para lhe orientar a respeito de como incluir novos
% itens. Obs. os itens já listados podem ser apagados.
%
% De forma a não danificar a estruturação principal, as bibliotecas individuais
% utilizadas podem ser adicionadas no arquivo my_packages.tex localizado na pasta
% setup.
%
% Ao longo do documento você encontrará exemplos de como adicionar uma figura, uma
% tabela, fazer uma citação direta (longa e curta) e indireta.
%
% As listas e sumário são atualizados automaticamente.
% 
% Suas referencias devem ser colocadas no formato bibtex no arquivo references.bib
% na pasta raiz.
%
%                               Atualizações
%
%  04-03 - incorporada o pacote listings para inclusão
%          de arquivos de programação código c ou outras
%          linguagens no trabalho.
%
%  27-03 - ajuste do tamanho da fonte, alinhamento e 
%          posicionamento do caption da imagem
%
%  21-06 - Impressão da lista de simbolos
%
% Atenciosamente 
%
% Luiz Fernando M. Arruda - 21/06/2022

% ------------------------------------------------------------------------
% ------------------------------------------------------------------------

\documentclass[% -- opções da classe uniavan --
	12pt,				% tamanho da fonte
	openright,			% capítulos começam em pág ímpar (insere página vazia caso preciso)
	oneside,			% para impressão no anverso. Oposto a twoside
	% -- opções da classe abntex2 --
	chapter=TITLE,		% títulos de capítulos convertidos em letras maiúsculas
	section=TITLE,		% títulos de seções convertidos em letras maiúsculas
	%subsection=TITLE,	% títulos de subseções convertidos em letras maiúsculas
	%subsubsection=TITLE,% títulos de subsubseções convertidos em letras maiúsculas
	% -- opções do pacote babel --
	english,			% idioma adicional para hifenização
	%french,			% idioma adicional para hifenização
	%spanish,			% idioma adicional para hifenização
	brazil				% o último idioma é o principal do documento
	]{setup/uniavan}

\addbibresource{references.bib} % Seus arquivos de referências

% -------------------------------------------------------------------------
% Inclusão de novas bibliotecas
% -------------------------------------------------------------------------
% Inclusão de novas bibliotecas devem ser feitas arquivo setup/my_packages.tex
% -----------------------------

\input{setup/my_packages}

% -------------------------------------------------------------------------
% Dados gerais do documento
% -------------------------------------------------------------------------
% Inclusão do nome do autor, título, orientador, membros da banca, curso etc
% Devem ser editados no arquivo setup/information.tex
% -----------------------------

\input{setup/information}
% ---

% ---
% Configurações de aparência do PDF final
% ---
% alterando o aspecto da cor azul
\definecolor{blue}{RGB}{41,5,195}
% informações do PDF
\makeatletter
\hypersetup{
     	%pagebackref=true,
		pdftitle={\@title}, 
		pdfauthor={\@author},
    	pdfsubject={\imprimirpreambulo},
	    pdfcreator={LaTeX with abnTeX2},
		pdfkeywords={ufsc, latex, abntex2}, 
		colorlinks=true,       		% false: boxed links; true: colored links
    	linkcolor=black,%blue,          	% color of internal links
    	citecolor=black,%blue,        		% color of links to bibliography
    	filecolor=black,%magenta,      		% color of file links
		urlcolor=black,%blue,
		bookmarksdepth=4
}
\makeatother
% ---

% ---
% carregar e compilar a lista de abreviaturas e siglas e a lista de símbolos
% ---

% carrega a lista de abreviaturas e siglas e a lista de símbolos
\input{setup/list_of_acron_symbols}

% compila a lista de abreviaturas e siglas e a lista de símbolos
\makenoidxglossaries 

% ---

% ---
% compila o indice
% ---
\makeindex
% ---

% ----
% Início do documento
% ----
\begin{document}

% Seleciona o idioma do documento (conforme pacotes do babel)
%\selectlanguage{english}
\selectlanguage{brazil}

% Retira espaço extra obsoleto entre as frases.
\frenchspacing 

% Espaçamento 1.5 entre linhas
\OnehalfSpacing

% Corrige justificação
%\sloppy

% -------------------------------------------------------------------------
% ELEMENTOS PRÉ-TEXTUAIS
% -------------------------------------------------------------------------
% Capa, folha de rosto, ficha bibliográfica, errata, folha de aprovação
% Dedicatória, agradecimentos, epígrafe, resumos, listas
% -------------------------------------------------------------------------
% Capa
% -------------------------------------------------------------------------
\imprimircapa
% -------------------------------------------------------------------------


% -------------------------------------------------------------------------
% Folha de rosto
% (o * indica que haverá a ficha bibliográfica)
% -------------------------------------------------------------------------
\imprimirfolhaderosto
% -------------------------------------------------------------------------

% -------------------------------------------------------------------------
% Folha de aprovação
% -------------------------------------------------------------------------
% A folha de aprovação é feita em arquivo separado (para facilitar a impressão / eventual mudança de membro pelo suplente / inclusão da folha escaneada).
% O arquivo de edição da folha de aprovação  está em beforetext/aprovacao.tex
% -------------------------------------------------------------------------
%\input{beforetext/aprovacao}
%\includepdf{beforetext/aprovacao.pdf}
% -------------------------------------------------------------------------


% -------------------------------------------------------------------------
% Elementos pré-textuais não obrigatórios
% -------------------------------------------------------------------------
% Os elementos pré-textuais devem ser editados na pasta beforetext
% -------------------------------------------------------------------------
% Agradecimentos
% -------------------------------------------------------------------------
%\input{beforetext/agradecimentos}
% -------------------------------------------------------------------------


% -------------------------------------------------------------------------
% Epígrafe
% -------------------------------------------------------------------------
%\input{beforetext/epigrafe}
% -------------------------------------------------------------------------

% -------------------------------------------------------------------------
% Resumo
% -------------------------------------------------------------------------
\input{beforetext/resumo}
% -------------------------------------------------------------------------


% -------------------------------------------------------------------------
% Abstract
% -------------------------------------------------------------------------
\input{beforetext/abstract}
% -------------------------------------------------------------------------

% -------------------------------------------------------------------------
% LISTAS (não obrigatórias)
% -------------------------------------------------------------------------

% -------------------------------------------------------------------------
% Lista de figuras
% -------------------------------------------------------------------------
\imprimrilistadefiguras
% -------------------------------------------------------------------------
	
% -------------------------------------------------------------------------
% Lista de tabelas
% -------------------------------------------------------------------------
\imprimirlistadetabelas
% -------------------------------------------------------------------------

% -------------------------------------------------------------------------
% Lista de abreviaturas e siglas (devem ser declarados no preambulo)
% -------------------------------------------------------------------------
\imprimirlistadesiglas
% -------------------------------------------------------------------------

% -------------------------------------------------------------------------
% Lista de simbolos (devem ser declarados no preambulo)
% -------------------------------------------------------------------------
\imprimirlistadesimbolos
% -------------------------------------------------------------------------

% -------------------------------------------------------------------------
% Sumário
% -------------------------------------------------------------------------
\imprimirsumario
% -------------------------------------------------------------------------	

% ----------------------------------------------------------
% ELEMENTOS TEXTUAIS
% ----------------------------------------------------------
\textual

% ---

% ---
\input{chapters/01-chapter}
\input{chapters/02-chapter}
\input{chapters/03-chapter}
\input{chapters/04-chapter}
\input{chapters/05-chapter}

% ----------------------------------------------------------
% ELEMENTOS PÓS-TEXTUAIS
% ----------------------------------------------------------
\postextual
% ----------------------------------------------------------

% ----------------------------------------------------------
% Referências bibliográficas
% ----------------------------------------------------------
\begingroup
    \printbibliography[title=REFERÊNCIAS]
\endgroup

% ----------------------------------------------------------
% Glossário
% ----------------------------------------------------------
\imprimirglossario

% ----------------------------------------------------------
% Apêndices
% ----------------------------------------------------------

% ---
% Inicia os apêndices
% ---
\begin{apendicesenv}
	%\partapendices* 
	\input{aftertext/apendice_a}
	\input{aftertext/apendice_b}
\end{apendicesenv}
% ---


% ----------------------------------------------------------
% Anexos
% ----------------------------------------------------------

% ---
% Inicia os anexos
% ---
\begin{anexosenv}
	%\partanexos*
	\input{aftertext/anexo_a}
\end{anexosenv}

%---------------------------------------------------------------------
% INDICE REMISSIVO
%---------------------------------------------------------------------
% \phantompart
% \printindex
%---------------------------------------------------------------------

\end{document}
